% Options for packages loaded elsewhere
\PassOptionsToPackage{unicode}{hyperref}
\PassOptionsToPackage{hyphens}{url}
\PassOptionsToPackage{table}{xcolor}
%
\documentclass[
  letterpaper,
]{article}
\usepackage{amsmath,amssymb}
\usepackage{iftex}
\usepackage[T1]{fontenc}
\usepackage[utf8]{inputenc}
\usepackage{textcomp}
\usepackage{lmodern}
\ifPDFTeX\else
  % xetex/luatex font selection
\fi
% Use upquote if available, for straight quotes in verbatim environments
\IfFileExists{upquote.sty}{\usepackage{upquote}}{}
\IfFileExists{microtype.sty}{% use microtype if available
  \usepackage[]{microtype}
  \UseMicrotypeSet[protrusion]{basicmath} % disable protrusion for tt fonts
}{}
\makeatletter
\@ifundefined{KOMAClassName}{% if non-KOMA class
  \IfFileExists{parskip.sty}{%
    \usepackage{parskip}
  }{% else
    \setlength{\parindent}{0pt}
    \setlength{\parskip}{6pt plus 2pt minus 1pt}}
}{% if KOMA class
  \KOMAoptions{parskip=half}}
\makeatother
\usepackage{xcolor}
\usepackage{color}
\usepackage{fancyvrb}
\newcommand{\VerbBar}{|}
\newcommand{\VERB}{\Verb[commandchars=\\\{\}]}
\DefineVerbatimEnvironment{Highlighting}{Verbatim}{commandchars=\\\{\}}
% Add ',fontsize=\small' for more characters per line
\newenvironment{Shaded}{}{}
\newcommand{\AlertTok}[1]{\textcolor[rgb]{1.00,0.00,0.00}{\textbf{#1}}}
\newcommand{\AnnotationTok}[1]{\textcolor[rgb]{0.38,0.63,0.69}{\textbf{\textit{#1}}}}
\newcommand{\AttributeTok}[1]{\textcolor[rgb]{0.49,0.56,0.16}{#1}}
\newcommand{\BaseNTok}[1]{\textcolor[rgb]{0.25,0.63,0.44}{#1}}
\newcommand{\BuiltInTok}[1]{\textcolor[rgb]{0.00,0.50,0.00}{#1}}
\newcommand{\CharTok}[1]{\textcolor[rgb]{0.25,0.44,0.63}{#1}}
\newcommand{\CommentTok}[1]{\textcolor[rgb]{0.38,0.63,0.69}{\textit{#1}}}
\newcommand{\CommentVarTok}[1]{\textcolor[rgb]{0.38,0.63,0.69}{\textbf{\textit{#1}}}}
\newcommand{\ConstantTok}[1]{\textcolor[rgb]{0.53,0.00,0.00}{#1}}
\newcommand{\ControlFlowTok}[1]{\textcolor[rgb]{0.00,0.44,0.13}{\textbf{#1}}}
\newcommand{\DataTypeTok}[1]{\textcolor[rgb]{0.56,0.13,0.00}{#1}}
\newcommand{\DecValTok}[1]{\textcolor[rgb]{0.25,0.63,0.44}{#1}}
\newcommand{\DocumentationTok}[1]{\textcolor[rgb]{0.73,0.13,0.13}{\textit{#1}}}
\newcommand{\ErrorTok}[1]{\textcolor[rgb]{1.00,0.00,0.00}{\textbf{#1}}}
\newcommand{\ExtensionTok}[1]{#1}
\newcommand{\FloatTok}[1]{\textcolor[rgb]{0.25,0.63,0.44}{#1}}
\newcommand{\FunctionTok}[1]{\textcolor[rgb]{0.02,0.16,0.49}{#1}}
\newcommand{\ImportTok}[1]{\textcolor[rgb]{0.00,0.50,0.00}{\textbf{#1}}}
\newcommand{\InformationTok}[1]{\textcolor[rgb]{0.38,0.63,0.69}{\textbf{\textit{#1}}}}
\newcommand{\KeywordTok}[1]{\textcolor[rgb]{0.00,0.44,0.13}{\textbf{#1}}}
\newcommand{\NormalTok}[1]{#1}
\newcommand{\OperatorTok}[1]{\textcolor[rgb]{0.40,0.40,0.40}{#1}}
\newcommand{\OtherTok}[1]{\textcolor[rgb]{0.00,0.44,0.13}{#1}}
\newcommand{\PreprocessorTok}[1]{\textcolor[rgb]{0.74,0.48,0.00}{#1}}
\newcommand{\RegionMarkerTok}[1]{#1}
\newcommand{\SpecialCharTok}[1]{\textcolor[rgb]{0.25,0.44,0.63}{#1}}
\newcommand{\SpecialStringTok}[1]{\textcolor[rgb]{0.73,0.40,0.53}{#1}}
\newcommand{\StringTok}[1]{\textcolor[rgb]{0.25,0.44,0.63}{#1}}
\newcommand{\VariableTok}[1]{\textcolor[rgb]{0.10,0.09,0.49}{#1}}
\newcommand{\VerbatimStringTok}[1]{\textcolor[rgb]{0.25,0.44,0.63}{#1}}
\newcommand{\WarningTok}[1]{\textcolor[rgb]{0.38,0.63,0.69}{\textbf{\textit{#1}}}}
\usepackage{longtable,booktabs,array}
\usepackage{calc} % for calculating minipage widths
% Correct order of tables after \paragraph or \subparagraph
\usepackage{etoolbox}
\makeatletter
\patchcmd\longtable{\par}{\if@noskipsec\mbox{}\fi\par}{}{}
\makeatother
% Allow footnotes in longtable head/foot
\IfFileExists{footnotehyper.sty}{\usepackage{footnotehyper}}{\usepackage{footnote}}
\makesavenoteenv{longtable}
\setlength{\emergencystretch}{3em} % prevent overfull lines
\providecommand{\tightlist}{%
  \setlength{\itemsep}{0pt}\setlength{\parskip}{0pt}}
\setcounter{secnumdepth}{5}
% Self-contained presentation rules embedded in every generated LaTeX file.
\usepackage{iftex}
\ifPDFTeX\else\PackageError{SSE-ASAS}{Compile with pdfLaTeX only}{}\fi
\usepackage[T1]{fontenc}
\usepackage[utf8]{inputenc}
\usepackage{lmodern}
\usepackage{microtype}
\usepackage{geometry}
\geometry{margin=0.93in,headheight=16pt}
\usepackage{xcolor} % [table] passed at top of file (option-clash fix)
\definecolor{ink}{HTML}{15283B}
\definecolor{brand}{HTML}{075B75}
\definecolor{brandlight}{HTML}{EAF5F8}
\definecolor{warm}{HTML}{A15C1D}
\definecolor{muted}{HTML}{536577}
\usepackage{fancyhdr}
\usepackage{titlesec}
\usepackage{enumitem}
\usepackage{booktabs}
\usepackage{longtable}
\usepackage{array}
\usepackage{tabularx}
\usepackage{tcolorbox}
\tcbuselibrary{breakable,skins}
\usepackage{listings}
\usepackage{fancyvrb}
\usepackage{hyperref}
\hypersetup{colorlinks=true,linkcolor=brand,urlcolor=brand,citecolor=brand}
\urlstyle{same}
\usepackage{lastpage}
\setlength{\emergencystretch}{3em}
\setlength{\parskip}{0.45em}
\setlength{\parindent}{0pt}
\setlist[itemize]{leftmargin=*,topsep=3pt,itemsep=2pt}
\setlist[enumerate]{leftmargin=*,topsep=3pt,itemsep=2pt}
\titleformat{\section}{\Large\bfseries\color{ink}}{\thesection}{0.75em}{}[\vspace{-2pt}\color{brand}\titlerule]
\titleformat{\subsection}{\large\bfseries\color{brand}}{\thesubsection}{0.7em}{}
\titleformat{\subsubsection}{\normalsize\bfseries\color{ink}}{\thesubsubsection}{0.7em}{}
\pagestyle{fancy}
\fancyhf{}
\fancyhead[L]{\small\sffamily\textcolor{brand}{SOFTWARE SECURITY ENGINEERING}}
\fancyhead[R]{\small\sffamily\textcolor{muted}{TECHNICAL SERIES}}
\fancyfoot[L]{\scriptsize\sffamily Version 1.0 \textbullet\ 14 October 2026}
\fancyfoot[R]{\scriptsize\sffamily Page \thepage\ of \pageref{LastPage}}
\renewcommand{\headrulewidth}{0.5pt}
\renewcommand{\footrulewidth}{0.2pt}
\setcounter{tocdepth}{2}
\setcounter{secnumdepth}{3}
\lstdefinestyle{sse}{basicstyle=\ttfamily\footnotesize,frame=single,framerule=0.4pt,rulecolor=\color{brand!40},backgroundcolor=\color{brandlight!45},numbers=left,numberstyle=\scriptsize\color{muted},numbersep=6pt,breaklines=true,breakatwhitespace=false,tabsize=2,showstringspaces=false,columns=fullflexible,keepspaces=true,captionpos=b,aboveskip=9pt,belowskip=9pt}
\lstset{style=sse}
\newtcolorbox{sourceconcept}[1][]{enhanced,breakable,colback=brandlight,colframe=brand!65!black,title={ORIGINAL CONCEPT},fonttitle=\bfseries\sffamily,#1}
\newtcolorbox{modernextension}[1][]{enhanced,breakable,colback=white,colframe=warm!85!black,title={MODERN EXTENSION},fonttitle=\bfseries\sffamily,#1}
\newtcolorbox{fieldwork}[1][]{enhanced,breakable,colback=white,colframe=ink!55,title={PRACTICAL APPLICATION},fonttitle=\bfseries\sffamily,#1}
\AtBeginDocument{\color{ink}}
\makeatletter
\renewcommand{\maketitle}{%
 \begin{titlepage}\thispagestyle{empty}
 \vspace*{1.5cm}{\color{brand}\rule{\textwidth}{3pt}}\par\vspace{1cm}
 {\sffamily\bfseries\color{muted}\large SOFTWARE SECURITY ENGINEERING\\ APPLICATION SECURITY ASSESSMENT SERIES}\par\vspace{1.05cm}
 {\sffamily\bfseries\color{ink}\fontsize{27}{34}\selectfont\@title\par}\vspace{0.8cm}
 {\sffamily\large\color{brand}\@author\par}\vfill
 \begin{tcolorbox}[colback=brandlight,colframe=brand!55,boxrule=0.5pt]
 \small\textbf{Editorial policy}\quad Original teaching synthesis with book-grounded concepts distinguished from contemporary practices. Modern platform behavior is version-sensitive; validate against linked vendor documentation before production use.
 \end{tcolorbox}
 {\color{brand}\rule{\textwidth}{2pt}}\par\vspace{0.3cm}
 {\sffamily\small Version 1.0 \quad|\quad October 2026 \quad|\quad pdfLaTeX}\end{titlepage}\setcounter{page}{1}}
\makeatother
\ifLuaTeX
  \usepackage{selnolig}  % disable illegal ligatures
\fi
\usepackage{bookmark}
\IfFileExists{xurl.sty}{\usepackage{xurl}}{} % add URL line breaks if available
\urlstyle{same}
\hypersetup{
  pdftitle={Static Analysis Algorithms and CodeQL},
  pdfauthor={Independent Technical Study Edition},
  hidelinks,
  pdfcreator={LaTeX via pandoc}}

\title{Static Analysis Algorithms and CodeQL}
\author{Independent Technical Study Edition}
\date{}

\begin{document}
\maketitle

{
\setcounter{tocdepth}{3}
\tableofcontents
}
\section{Objectives and provenance}\label{objectives-and-provenance}

This document explains enough static-analysis theory to interpret tool
results, reason about false negatives, write a small analyzer, and
design evidence-based triage. \textbf{Core source:} Chess and West,
Chapters 2--4, specifically parsing, abstract syntax, control-flow and
data-flow, taint propagation, pointer aliasing,
intraprocedural/interprocedural analysis, rules, and reporting.
\textbf{Modern extension:} CodeQL database/queries, code scanning
workflows, SARIF results, and GitHub Code Security.

\section{What static analysis can and cannot
establish}\label{what-static-analysis-can-and-cannot-establish}

Static analysis reasons about program behavior without executing every
possible input. It can enforce syntactic conventions, prove certain
properties under explicit assumptions, and flag potentially unsafe
execution paths. Soundness and completeness depend on the model,
language semantics, path abstraction, environmental assumptions, and
whether the tool aims for few false positives or wide coverage.

A \emph{sound over-approximation} for a property includes all relevant
runtime behaviors but may report infeasible ones. A deliberately narrow,
high-precision query may miss vulnerable flows, especially through
reflection, native interfaces, unmodeled frameworks, dynamic dispatch,
or unusual data sources. Interpret ``no findings'' as ``no findings
under the configured analysis,'' never as a proof of safety.

\section{Front end and program
representation}\label{front-end-and-program-representation}

\subsection{Lexing, parsing, types, and abstract
syntax}\label{lexing-parsing-types-and-abstract-syntax}

A typical analyzer tokenizes source, constructs an abstract syntax tree
(AST), resolves names, checks or approximates semantics, and builds
control/data-flow representations. A lexical search can find
\texttt{system(...)}; an AST allows distinguishing function calls from
comments and string literals. Call graphs approximate interprocedural
relationships; control-flow graphs (CFGs) represent possible execution
paths.

\textbf{Example CFG:} for
\texttt{if\ trusted:\ x\ =\ input;\ else:\ x\ =\ "";\ use(x)}, the CFG
branches on \texttt{trusted} and merges before \texttt{use}. Reaching
definitions and taint facts at the merge determine whether
attacker-derived data can reach \texttt{use}. Even if \texttt{trusted}
is always false in deployment, a path-insensitive analysis may keep the
positive branch.

\subsection{Static single assignment and phi
merging}\label{static-single-assignment-and-phi-merging}

SSA uses one definition per variable version. For a branch writing
\texttt{x1} or \texttt{x2}, the merge introduces \(x_3=\phi(x_1,x_2)\).
That makes def-use relationships explicit. Treat phi functions as
analysis representations, not executable instructions. Chess and West
give these concepts and examples in Chapter 4.

\section{Data-flow equations and
fixpoint}\label{data-flow-equations-and-fixpoint}

Let \(CFG=(N,E)\). Associate each node \(n\) with facts \(IN_n\) and
\(OUT_n\) drawn from finite lattice \(L\). For a forward may-analysis:

\[IN_n=\bigsqcup_{p\in pred(n)} OUT_p,\qquad OUT_n=F_n(IN_n).\]

For a powerset taint domain, join is set union. Each transfer \(F_n\)
propagates or removes facts according to the modeled operation. Iterate
until a fixpoint is reached. For a finite-height lattice and monotone
transfer functions, a worklist algorithm terminates.

\begin{Shaded}
\begin{Highlighting}[]
\NormalTok{worklist := all CFG nodes}
\NormalTok{IN[n] := bottom; OUT[n] := bottom}
\NormalTok{while worklist not empty:}
\NormalTok{    n := take(worklist)}
\NormalTok{    merged := join(OUT[p] for p in pred[n])}
\NormalTok{    changed := (transfer(n, merged) != OUT[n])}
\NormalTok{    IN[n] := merged}
\NormalTok{    OUT[n] := transfer(n, merged)}
\NormalTok{    if changed: add all successors(n) to worklist}
\end{Highlighting}
\end{Shaded}

In a may-taint analysis, \texttt{tainted} at a sink means there is a
modeled path \emph{possibly} conveying untrusted data. Feasibility of
the path must still be established. For must-analysis or different
lattices, join and initialization differ; do not reuse a union-based
algorithm without checking semantics.

\section{Taint rules: sources, transfers, sanitizers,
sinks}\label{taint-rules-sources-transfers-sanitizers-sinks}

\begin{sourceconcept}
Chess and West distinguish **sources**, **sinks**, **pass-through rules**, **cleanse rules**, and **entry-point rules**. A flow is reportable when attacker-controlled data can reach a sensitive operation under the analysis model without an appropriate contextual defense.
\end{sourceconcept}

Consider the path
\texttt{request\ parameter\ -\textgreater{}\ trim()\ -\textgreater{}\ concatenate\ SQL\ -\textgreater{}\ execute()}.
A \texttt{trim()} call propagates taint. Parameter binding changes the
interpretation boundary and prevents data from becoming SQL grammar,
whereas trimming whitespace does not cleanse SQL taint.

\begin{Shaded}
\begin{Highlighting}[]
\NormalTok{user }\OperatorTok{=}\NormalTok{ request.args[}\StringTok{\textquotesingle{}name\textquotesingle{}}\NormalTok{]   }\CommentTok{\# source}
\NormalTok{user }\OperatorTok{=}\NormalTok{ user.strip()           }\CommentTok{\# propagator}
\NormalTok{query }\OperatorTok{=} \StringTok{"SELECT * FROM staff WHERE name=\textquotesingle{}"} \OperatorTok{+}\NormalTok{ user }\OperatorTok{+} \StringTok{"\textquotesingle{}"}
\NormalTok{cursor.execute(query)         }\CommentTok{\# SQL sink}
\end{Highlighting}
\end{Shaded}

A field-specific sanitizer must be valid for the \emph{particular sink
context}. Treating generic HTML escaping as cleansing an SQL query is a
modeling bug. Sanitization is not merely removing ``bad characters''; it
can be a type/range proof, safe API, parameter binding, allowlisted
token, or context-specific transformation.

\subsection{Sensitivity and precision
trade-offs}\label{sensitivity-and-precision-trade-offs}

\begin{longtable}[]{@{}lll@{}}
\toprule\noalign{}
Dimension & Insensitive approach & Sensitive approach \\
\midrule\noalign{}
\endhead
\bottomrule\noalign{}
\endlastfoot
Flow & Ignores program order & Uses reaching definitions/order \\
Context & Merges all call sites & Distinguishes callers/call strings \\
Field & Marks whole object & Tracks individual fields \\
Path & Joins both branches & Reasons about predicates \\
Alias & Treats pointers broadly & Models may-/must-alias facts \\
\end{longtable}

Greater sensitivity can reduce false positives but increases cost and
complexity. A tool can still miss flows through unknown libraries or
incomplete extraction. Record precision and coverage trade-offs whenever
selecting query suites.

\section{Worked example: minimal taint
trace}\label{worked-example-minimal-taint-trace}

\begin{Shaded}
\begin{Highlighting}[]
\KeywordTok{def}\NormalTok{ route(req):}
\NormalTok{    raw }\OperatorTok{=}\NormalTok{ req.query[}\StringTok{\textquotesingle{}filename\textquotesingle{}}\NormalTok{]}
\NormalTok{    safe }\OperatorTok{=}\NormalTok{ validate\_slug(raw)}
    \ControlFlowTok{return}\NormalTok{ storage.}\BuiltInTok{open}\NormalTok{(safe)}
\end{Highlighting}
\end{Shaded}

\textbf{Trace step 1:} entry-point modeling determines whether
\texttt{req.query} produces attacker-controlled values. \textbf{Step 2:}
data flow follows \texttt{raw} to \texttt{validate\_slug}; its
implementation and preconditions determine whether it can be modeled as
a sanitizer. \textbf{Step 3:} \texttt{storage.open} is a sensitive path
sink only if file selection crosses an authorization or filesystem
boundary. \textbf{Step 4:} context asks whether the storage API already
confines names to a fixed namespace. Without source/sink library models,
the analyzer can miss this path completely.

\section{CodeQL and modern SAST
operations}\label{codeql-and-modern-sast-operations}

\subsection{CodeQL as a database and query
pipeline}\label{codeql-as-a-database-and-query-pipeline}

CodeQL represents source in a per-language database and evaluates
queries over it. In advanced cases, custom queries identify
domain-specific sources and sinks and model libraries. Start with
official built-in query suites, inspect their rule coverage, and extend
where the organization's frameworks are not modeled. Query packs
distribute rule logic; library/model packs carry reusable extraction and
modeling logic where supported.

\textbf{CLI illustration} (illustrative Python project; verify supported
CLI and pack versions):

\begin{Shaded}
\begin{Highlighting}[]
\ExtensionTok{codeql}\NormalTok{ database create db{-}python }\AttributeTok{{-}{-}language}\OperatorTok{=}\NormalTok{python }\AttributeTok{{-}{-}source{-}root}\OperatorTok{=}\NormalTok{.}
\ExtensionTok{codeql}\NormalTok{ database analyze db{-}python codeql/python{-}queries }\DataTypeTok{\textbackslash{}}
  \AttributeTok{{-}{-}format}\OperatorTok{=}\NormalTok{sarif{-}latest }\AttributeTok{{-}{-}output}\OperatorTok{=}\NormalTok{results.sarif}
\end{Highlighting}
\end{Shaded}

\textbf{Default versus advanced setup:} GitHub recommends default setup
as the low-maintenance starting point. Advanced setup gives explicit
workflow/build/query configuration when necessary. GitHub documents
distinct \texttt{none}, \texttt{autobuild}, and \texttt{manual} build
modes for compiled languages; choose based on actual language and CI
support rather than assuming a single build recipe.

\subsection{Pull-request and branch
semantics}\label{pull-request-and-branch-semantics}

Scanning default branch periodically is insufficient when a
branch-protection rule requires evidence \emph{before merging}. Ensure
that scans execute on pull requests and required check contexts refer to
those runs. Verify supported commit/branch association, relevant
protected branches, permissions, forks and trust model, and action
versions. Test a deliberately seeded, harmless lab finding to establish
that the intended rule blocks only when the expected check fails.

\begin{modernextension}
GitHub presently documents `default` and `security-extended` CodeQL suites. The latter adds broader coverage at potentially lower precision. GitHub Advanced Security is the umbrella term, while **GitHub Code Security** and **GitHub Secret Protection** are distinct product entitlements for certain private/internal repository features. Confirm availability for the repository and enterprise plan.
\end{modernextension}

\section{SARIF, triage, and
verification}\label{sarif-triage-and-verification}

SARIF 2.1.0 carries tool identity, rule IDs, result locations, severity
metadata, and, when produced by an analyzer, code-flow paths. Alert
state and confidence are different from exploitability. For each alert:
establish relevant input, check actual sink, verify framework defaults,
evaluate required privilege and reachability, assess mitigating
controls, identify fix and test, and retain uncertainty.

\begin{longtable}[]{@{}
  >{\raggedright\arraybackslash}p{(\columnwidth - 4\tabcolsep) * \real{0.3333}}
  >{\raggedright\arraybackslash}p{(\columnwidth - 4\tabcolsep) * \real{0.3333}}
  >{\raggedright\arraybackslash}p{(\columnwidth - 4\tabcolsep) * \real{0.3333}}@{}}
\toprule\noalign{}
\begin{minipage}[b]{\linewidth}\raggedright
Signal
\end{minipage} & \begin{minipage}[b]{\linewidth}\raggedright
What it shows
\end{minipage} & \begin{minipage}[b]{\linewidth}\raggedright
What it cannot prove
\end{minipage} \\
\midrule\noalign{}
\endhead
\bottomrule\noalign{}
\endlastfoot
Static path & Modeled data relation & Runtime feasibility without more
evidence \\
Alert severity & Rule/tool risk estimate & Product-specific business
impact \\
Clean scan & No matching findings & Complete code coverage or absence of
flaws \\
Fixed alert & Analyzer no longer reports it on analyzed ref & Full
security assurance without testing \\
CodeQL model extension & New modeled source/sink behavior & Automatic
correctness of all custom models \\
\end{longtable}

\textbf{Measures:} scan coverage by active repository and language;
successful scan rate; median time to triage; confirmed vulnerability
rate; age of open high-risk findings; number of approved and monitored
exceptions; test pass rate on fixes. Avoid using raw alert counts as
engineer productivity scores.

\section{Exercises and acceptance
evidence}\label{exercises-and-acceptance-evidence}

\begin{enumerate}
\def\labelenumi{\arabic{enumi}.}
\tightlist
\item
  Draw CFG and compute taint sets for a branching string-handling
  routine.
\item
  Mark source, propagator, source-to-sink path, contextual sanitizer,
  and sink in a short query.
\item
  Modify a synthetic example so an apparent path is infeasible; describe
  whether a path-insensitive analyzer still reports it.
\item
  Create a CodeQL database for an authorized practice repository and
  explain extraction coverage.
\item
  Document one real code scanning alert with a validated trace, a
  targeted fix, and a regression test.
\end{enumerate}

\textbf{Exit criteria:} show a correct worklist trace, explain at least
three sensitivity dimensions, reproduce a SARIF result, document a
negative test, and articulate one confirmed limitation of the selected
analyzer.

\section{Annex A: a worked forward data-flow
solution}\label{annex-a-a-worked-forward-data-flow-solution}

Consider a small fictional program with basic blocks \(b_0\) (entry),
\(b_1\) (trusted assignment), \(b_2\) (untrusted assignment), \(b_3\)
(merge), \(b_4\) (query). The entry branch may execute either \(b_1\) or
\(b_2\); both reach \(b_3\). Let the abstract fact set be either
\(\varnothing\) (the variable is not known to be tainted on any modeled
path) or \(\{x\}\) (taint may reach the variable). At entry, use
\(IN_{b_0}=\varnothing\). The trusted assignment kills taint, so
\(OUT_{b_1}=\varnothing\). The untrusted assignment generates taint, so
\(OUT_{b_2}=\{x\}\). At merge,
\(IN_{b_3}=OUT_{b_1}\cup OUT_{b_2}=\{x\}\), so the sink at \(b_4\) is
reported by a may-taint policy. The observation proves a path \emph{may}
be modeled; it does not prove the branch is feasible for an actual
attacker.

\textbf{Monotonicity:} as an analysis discovers additional possible
tainted paths, previously discovered taint facts cannot vanish solely
because the worklist visits a node in a different order. For finite
sets, union and monotone transfer functions guarantee convergence.
Flow-sensitive \emph{kill} operations must be modeled as part of a
transfer; they do not invalidate the monotonic framework if defined
correctly. In a large analyzer, widening may be needed for more
expressive infinite-height abstractions.

\subsection{Taint is a type of information-flow
question}\label{taint-is-a-type-of-information-flow-question}

A binary taint flag is useful but often too coarse. Input may be safe
for SQL when bound as a value, yet unsafe as a filesystem path. A
practical analyzer can track categories like \texttt{SQL\_VALUE},
\texttt{SHELL\_COMMAND}, \texttt{PATH\_COMPONENT}, and
\texttt{HTML\_MARKUP}, together with source provenance. Its domain and
transfer rules must form a coherent model; merely adding more string
flags without well-defined joins, sanitizers, and sink interpretations
creates inconsistent results.

For a simplified domain \(L=2^{\{sql,path,html\}}\), an operation can
remove one property while preserving others. An HTML-escaping function
should not erase the \texttt{sql} or \texttt{path} property. This gives
a principled explanation for context-sensitive cleansing and for why
broadly marking a value \texttt{safe} can produce severe false
negatives.

\subsection{Interprocedural summaries and framework
modeling}\label{interprocedural-summaries-and-framework-modeling}

A function summary records how outputs depend on inputs and whether it
creates, propagates, or cleanses taint. For example,
\texttt{normalize(String\ x)} may return data derived from \texttt{x},
while \texttt{executeBoundQuery(query,\ args)} may be safe with respect
to SQL syntax if its implementation correctly delegates argument
binding. Dynamic dispatch, reflection, dependency injection, and
generated code complicate identifying the function that actually
executes. If the analyzer assumes the wrong target, a complete
within-function dataflow is still insufficient.

A reviewer should therefore ask: does the CodeQL query understand our
HTTP framework's user-controlled entry points? Does it model the ORM
wrapper, custom validators, background events, and internal file
resolver? Does the build include generated source and all relevant
language units? The answers may require a model pack, a custom query, a
different build mode, or a separate manual review, depending on GitHub's
current support.

\section{Annex B: SAST operational decision
records}\label{annex-b-sast-operational-decision-records}

\begin{longtable}[]{@{}
  >{\raggedright\arraybackslash}p{(\columnwidth - 4\tabcolsep) * \real{0.3333}}
  >{\raggedright\arraybackslash}p{(\columnwidth - 4\tabcolsep) * \real{0.3333}}
  >{\raggedright\arraybackslash}p{(\columnwidth - 4\tabcolsep) * \real{0.3333}}@{}}
\toprule\noalign{}
\begin{minipage}[b]{\linewidth}\raggedright
Question
\end{minipage} & \begin{minipage}[b]{\linewidth}\raggedright
Required observation
\end{minipage} & \begin{minipage}[b]{\linewidth}\raggedright
Corrective action if missing
\end{minipage} \\
\midrule\noalign{}
\endhead
\bottomrule\noalign{}
\endlastfoot
Was the intended ref scanned? & SHA, branch/PR, successful run & Fix
event triggers and required checks \\
Was code extracted? & Supported languages, files and extraction logs &
Update build mode, dependencies, or matrix \\
Is rule coverage fit for purpose? & Query suite and included rules & Add
targeted suite/custom query after testing \\
Is alert relevant? & Modeled trace, actual call path, context & Manual
inspection and source/model correction \\
Is a fix validated? & Negative test plus scan of fix ref & Require
test/scan on PR before merge \\
Can results be correlated? & Stable SARIF category and tool ID &
Normalize uploads and branch association \\
\end{longtable}

\textbf{SARIF review notes:} tool output has locations and sometimes
rich paths; it is not a business impact analysis. Use stable identifiers
for rules, avoid changing SARIF run identity without reason, and verify
uploaded results belong to the correct revision. If a tool writes a
SARIF file but the upload action fails, the GitHub code scanning
interface has no validated new result from that step.

\subsection{A minimal investigation for a custom
query}\label{a-minimal-investigation-for-a-custom-query}

Define the target weakness, one unambiguous positive, and multiple
negatives. Specify sources and sinks in terms of framework semantics.
Develop the smallest query model that distinguishes positive from
negative; run against synthetic samples and a known repository subset.
Check whether a library call, sanitizer, and interprocedural flow are
captured. Review performance and false positives before enforcing the
query as a merge gate. Preserve query pack version and test fixtures in
the same change-controlled lifecycle as code.

\textbf{Limit statement to include in any analysis report:} ``Results
are conditioned on supported language extraction, the selected
query/rule packs, modeled source and sink APIs, scan event coverage, and
analyzer version. Untested paths and business-invariant violations may
remain.'' This is accurate even when the scan is clean.

\section{Source and modern extension
map}\label{source-and-modern-extension-map}

\textbf{Book-grounded:} Chess \& West Chapters 2--4 for the theory,
static analysis internals, rule types, review workflow, and reporting;
Chapters 5--12 for defect examples and analysis exercises.

\textbf{Modern extension:}
\url{https://docs.github.com/en/code-security/concepts/code-scanning/codeql/codeql-cli};
\url{https://docs.github.com/en/code-security/how-tos/find-and-fix-code-vulnerabilities/configure-code-scanning/configure-code-scanning};
\url{https://docs.github.com/en/code-security/how-tos/find-and-fix-code-vulnerabilities/configure-code-scanning/configuring-advanced-setup-for-code-scanning};
\url{https://docs.github.com/en/code-security/how-tos/find-and-fix-code-vulnerabilities/integrate-with-existing-tools/upload-sarif-file}.

\end{document}
